\unit{Thuyết trình 232 -- Các lược đồ Picard: các định lý tồn tại}
\sid{232}
\markboth{Thuyết trình 232}{Thuyết trình 232}
\src{143}
\seminar{Năm thứ 14\textsuperscript{}, 1961/62, số 232\textsuperscript{}\\Tháng Hai 1962}
\paperhead{KỸ THUẬT HẠ XUỐNG VÀ CÁC ĐỊNH LÝ TỒN TẠI TRONG HÌNH HỌC ĐẠI SỐ\\
V. CÁC LƯỢC ĐỒ PICARD: CÁC ĐỊNH LÝ TỒN TẠI}%
  {do Alexander GROTHENDIECK}

\fsection{1}{Các nhóm và hàm tử Picard tương đối}

Với mọi tiền lược đồ (tổng quát hơn, mọi không gian có vành) $X$, ta gọi
\emph{nhóm Picard} (tuyệt đối) của $X$, và ký hiệu là
$\operatorname{Pic}(X)$, nhóm các lớp đẳng cấu của các Môđun khả nghịch
(nghĩa là đẳng cấu địa phương với $\mathcal O_X$) trên $X$. Do đó ta có một
đẳng cấu chính tắc
\begin{equation}
\operatorname{Pic}(X)\xrightarrow{\sim}H^1(X,\mathcal O_X^*),
\tag{1.1}
\end{equation}
trong đó $\mathcal O_X^*$ chỉ bó các đơn vị của $\mathcal O_X$ (quả thật bó
này đồng nhất với bó các tự đẳng cấu của Môđun khả nghịch chuẩn
$\mathcal O_X$). Lưu ý rằng $X\longmapsto\operatorname{Pic}(X)$ hiển nhiên là
một hàm tử phản biến theo $X$, và đẳng cấu (1.1) có tính hàm tử.

Nếu $X$ là một tiền lược đồ trên một tiền lược đồ khác $S$, thì khi $S'$ biến
thiên trong phạm trù $(\mathrm{Sch})/S$ các tiền lược đồ trên $S$, từ điều trên
ta có hàm tử phản biến
$S'\longmapsto\operatorname{Pic}(X\times_S S')$. Hàm tử này hoàn toàn không
thể «biểu diễn được» ([4], II, A), vì do tồn tại những tự đẳng cấu không tầm
thường của các Môđun khả nghịch cần phân loại, hàm tử này không «có bản chất
địa phương» ([5], IV, 5.4). Vì vậy cần «làm cho nó thành địa phương» bằng cách,
nói chung, với mọi tiền lược đồ tương đối $X/S$, đưa vào một nhóm có bản chất
tương đối,
\begin{equation}
\operatorname{Pic}'(X/S)=H^0\bigl(S,R^1f_*(\mathcal O_X^*)\bigr),
\tag{1.2}
\end{equation}
(trong đó $f:X\longrightarrow S$ là cấu xạ cấu trúc) (so sánh [4], II,
C 3). Ở chỗ đã dẫn, nhóm này được gọi là nhóm Picard tương đối; ở đây nên gọi
nó là \emph{nhóm Picard tương đối hạn chế} của $X/S$, vì những lý do sẽ xuất
hiện sau đây. Khi $S'$ biến thiên trong $(\mathrm{Sch})/S$,
$S'\longmapsto\operatorname{Pic}'(X\times_S S'/S')$ là một hàm tử phản biến
theo $S'$, cũng được ký hiệu là $\operatorname{Pic}'_{X/S}$, do đó về cơ bản
được định nghĩa bởi công thức
\begin{equation}
\operatorname{Pic}'_{X/S}(S')=
\operatorname{Pic}'(X\times_S S'/S').
\tag{1.3}
\end{equation}
\src{144}
Hàm tử này nay «có bản chất địa phương», vì ta đã làm điều cần thiết cho mục
đích ấy. Một cách trực giác, vế phải của (1.3) được hiểu là tập các «họ đại số»
của những lớp bó khả nghịch trên (các thớ của) $X/S$, được đánh chỉ số bởi tiền
lược đồ tham số $S'/S$. Khi hàm tử $\operatorname{Pic}'$ biểu diễn được, tiền
lược đồ trên $S$ biểu diễn nó được ký hiệu là $\operatorname{Pic}'_{X/S}$ và
được gọi là \emph{tiền lược đồ Picard} của $X$ trên $S$; do đó khi ấy ta có
\begin{equation}
\operatorname{Hom}_S(S',\operatorname{Pic}'_{X/S})
\simeq\operatorname{Pic}'_{X/S}(S')
=\operatorname{Pic}'(X\times_S S'/S').
\tag{1.4}
\end{equation}

Tuy nhiên, có những trường hợp quan trọng trong đó
$\operatorname{Pic}'_{X/S}$ không biểu diễn được (ví dụ: một đa tạp
«Brauer--Severi» trên một trường $k$, không có điểm hữu tỉ trên $k$), nhưng vẫn
có một định nghĩa tự nhiên cho tiền lược đồ Picard tương đối. Điều này là do
trong định nghĩa hàm tử $\operatorname{Pic}'$ từ các nhóm Picard tuyệt đối
$\operatorname{Pic}(X\times_S S'/S')$, ta vẫn chưa địa phương hóa đủ; chính
xác hơn, nói chung $\operatorname{Pic}'$ không «tương thích với hạ phẳng trung
thành». Ta làm rõ điều này.

Gọi $\mathfrak M$ là lớp các cấu xạ tiền lược đồ \emph{phẳng trung thành và tựa
compact}; lớp này ổn định qua đổi cơ sở và qua hợp thành. Cho $P$ là một hàm tử
phản biến từ $(\mathrm{Sch})/S$ đến phạm trù các tập hợp, và với mọi
$S$-cấu xạ $u:T'\longrightarrow T$, $u\in\mathfrak M$, xét biểu đồ
\begin{equation}
P(T)\longrightarrow P(T')\rightrightarrows P(T'\times_TT')
\tag{1.5}
\end{equation}
là ảnh qua $P$ của biểu đồ
\[
T\longleftarrow T'\mathrel{\substack{\xleftarrow{\,\mathrm{pr}_1\,}\\[-0.35ex]
\xleftarrow[\,\mathrm{pr}_2\,]{}}}T'\times_TT'.
\]
Nếu $P$ biểu diễn được, thì từ lý thuyết hạ xuống ([4], I, B, định lý 2) suy ra
biểu đồ (1.5) là khớp với mọi $u\in\mathfrak M$. Ta diễn đạt sự kiện này bằng
cách nói $P$ tương thích với $\mathfrak M$, trong trường hợp này là $P$
«tương thích với hạ phẳng trung thành», hay «tiền bó» $P$ trên
$(\mathrm{Sch})/S$ là một «bó» đối với khái niệm địa phương hóa do lớp
$\mathfrak M$ cung cấp. Khi $P$ tùy ý, một quy trình chuẩn, quen thuộc trong
trường hợp địa phương hóa tôpô thông thường, cho phép liên kết với nó một «bó»
$\mathcal P$ và một đồng cấu hàm tử $P\longrightarrow\mathcal P$, phổ dụng
theo nghĩa hiển nhiên. Có thể tính $\mathcal P$ một cách tường minh như sau:
để định nghĩa $\mathcal P(T)$, với mọi $T'$ trên $T$ sao cho cấu xạ
$u:T'\longrightarrow T$ thuộc $\mathfrak M$, ta ký hiệu bởi
\src{145}
$\overline H^0(T'/T,P)$ tập con của $P(T')$ gồm các phần tử $\xi$ sao cho hai
ảnh $\xi_1,\xi_2$ của chúng trong $P(T'\times_TT')$ thỏa điều kiện: tồn tại
một cấu xạ $v:T''\longrightarrow T'\times_TT'$, $v\in\mathfrak M$, sao cho
$\xi_1$ và $\xi_2$ có cùng ảnh trong $P(T'')$. [Ghi chú. -- Tập
$\overline H^0$ được định nghĩa như vậy lớn hơn tập $H^0(T'/T,P)$ đã đưa vào
trong [4], I, A, 4(a).] Khi $T'$ biến thiên trên $T$ cố định (vẫn với
$u\in\mathfrak M$), các $\overline H^0(T'/T,P)$ tạo thành một hệ quy nạp (khi
tập các $T'$ được trang bị tiền thứ tự xác định bởi quan hệ trội), và ta đặt
\begin{equation}
\mathcal P(T)=\varinjlim_{T'}\overline H^0(T'/T,P).
\tag{1.6}
\end{equation}
Quy luật hàm tử theo $T$ của biểu thức này được tường minh hóa một cách hiển
nhiên.

Khi
\[
P(T)=\operatorname{Pic}(X\times_ST),
\]
hàm tử phản biến trên $(\mathrm{Sch})/S$ được định nghĩa bởi (1.6) được gọi là
\emph{hàm tử Picard tương đối} của $X$ trên $S$, và được ký hiệu là
$\operatorname{Pic}_{X/S}$; ta gọi \emph{nhóm Picard tương đối} của $X$ trên
$S$, và ký hiệu là $\operatorname{Pic}(X/S)$, nhóm
$\operatorname{Pic}_{X/S}(S)$. Khi đó ta có song ánh hiển nhiên
\begin{equation}
\operatorname{Pic}_{X/S}(T)\simeq\operatorname{Pic}(X\times_ST/T).
\tag{1.7}
\end{equation}

Do đó một phần tử của $\operatorname{Pic}(X/S)$ được định nghĩa nhờ một phần tử
$\xi'$ của một nhóm $\operatorname{Pic}(X\times_SS')$ (trong đó
$S'\longrightarrow S$ phẳng trung thành và tựa compact), sao cho có thể tìm
được một cấu xạ phẳng trung thành và tựa compact
$S''\longrightarrow S'\times_SS'$ sao cho hai ảnh ngược của $\xi'$ trong
$\operatorname{Pic}(X\times_SS'')$ bằng nhau. Một phần tử $\xi'$ của
$\operatorname{Pic}(X\times_SS')$ và $\xi_1$ của
$\operatorname{Pic}(X\times_SS_1)$ (thỏa các điều kiện vừa nêu) xác định cùng
một phần tử của $\operatorname{Pic}(X/S)$ khi và chỉ khi tồn tại một cấu xạ
phẳng trung thành và tựa compact $S'_1\longrightarrow S'\times_SS_1$ sao cho
ảnh của hai phần tử đang xét trong $\operatorname{Pic}(X\times_SS'_1)$ bằng
nhau. Thường vẫn tiện khi làm việc với hàm tử
$P'=\operatorname{Pic}'_{X/S}$ đã đưa vào ở trên; ta thấy ngay rằng cấu xạ
chính tắc $P\longrightarrow P'$ xác định một đẳng cấu
\begin{equation}
\mathcal P\xrightarrow{\sim}\mathcal P',
\tag{1.8}
\end{equation}
điều này cho một mô tả của $\operatorname{Pic}_{X/S}$ theo
$\operatorname{Pic}'_{X/S}=P'$, thường tiện hơn. Theo 2.3 dưới đây, khi trong
mô tả vừa nêu của $\operatorname{Pic}(X/S)$ ta thay $P$ bởi $P'$, ta thực ra
có thể
\src{146}
lấy $S''=S'\times_SS'$, $S'_1=S'\times_SS_1$, ít nhất dưới các điều kiện được
nêu rõ ở chỗ đã dẫn.

Khi hàm tử $\operatorname{Pic}_{X/S}$ biểu diễn được, ta nói $X/S$ có một tiền
lược đồ Picard, và tiền lược đồ trên $S$ biểu diễn hàm tử được gọi là
\emph{tiền lược đồ Picard} của $X$ trên $S$, vẫn được ký hiệu là
$\operatorname{Pic}_{X/S}$. Hiển nhiên, chỉ cần
$P'=\operatorname{Pic}'_{X/S}$ biểu diễn được, vì khi ấy $P'$ đã là một «bó»,
và công thức (1.8) chứng tỏ cấu xạ $P'\longrightarrow\mathcal P'$ đồng nhất
với cấu xạ chính tắc
\begin{equation}
\operatorname{Pic}'_{X/S}\longrightarrow\operatorname{Pic}_{X/S},
\tag{1.9}
\end{equation}
mà khi ấy là một đẳng cấu. Do đó thuật ngữ của ta tương thích với thuật ngữ đã
đưa vào ở trên cùng với (1.4). Nói chung, khi $\operatorname{Pic}_{X/S}$ tồn
tại, nó được định nghĩa bởi đẳng cấu có tính hàm tử
\begin{equation}
\operatorname{Hom}_S(S',\operatorname{Pic}_{X/S})\xrightarrow{\sim}
\operatorname{Pic}(X\times_SS'/S').
\tag{1.10}
\end{equation}

\fsection{2}{Quan hệ giữa các nhóm Picard tương đối và tuyệt đối khác nhau}

\result{MỆNH ĐỀ}{2.1} Cho $f:X\longrightarrow S$ là một cấu xạ sao cho
$\mathcal O_S\xrightarrow{\sim}f_*(\mathcal O_X)$. Khi đó ta có một dãy khớp
\[
0\longrightarrow\operatorname{Pic}(S)\longrightarrow\operatorname{Pic}(X)
\longrightarrow\operatorname{Pic}'(X/S).
\]
Khi $X$ có một tiết diện trên $S$, cấu xạ cuối là toàn ánh, nghĩa là ta có một
đẳng cấu
\[
\operatorname{Pic}'(X/S)\xrightarrow{\sim}
\operatorname{Pic}(X)/\operatorname{Pic}(S).
\]

Dãy khớp này có thể được xem là dãy khớp ở các bậc thấp tương ứng với dãy phổ
Leray cho $f$ và $\mathcal O_X^*$. Khẳng định thứ hai cũng có tính hình thức.

\result{MỆNH ĐỀ}{2.2} Cho $f:X\longrightarrow S$ là một cấu xạ tựa compact và
tách sao cho $\mathcal O_S\xrightarrow{\sim}f_*(\mathcal O_X)$, và cho
$S'\longrightarrow S$ là một cấu xạ phẳng trung thành và tựa compact. Khi đó:

\noindent (i) $\operatorname{Pic}'(X/S)\longrightarrow
\operatorname{Pic}'(X\times_SS'/S')$ là đơn ánh;

\noindent (ii) nếu $X$ có một tiết diện địa phương trên $S$ (nghĩa là mỗi
$s\in S$ có một lân cận mở $U$ sao cho $X|U$ có một tiết diện trên $U$), thì
biểu đồ
\[
\operatorname{Pic}'(X/S)\longrightarrow
\operatorname{Pic}'(X\times_SS'/S')\rightrightarrows
\operatorname{Pic}'(X\times_SS''/S'')
\]
\src{147}
(trong đó $S''=S'\times_SS'$) là khớp.

Khẳng định thứ nhất suy ra, nhờ các tính chất sơ cấp của hạ phẳng trung thành,
từ nhận xét tổng quát sau đây. Nếu $f:X\longrightarrow S$ là một cấu xạ sao cho
$\mathcal O_S\xrightarrow{\sim}f_*(\mathcal O_X)$, thì hàm tử
$\mathcal F\longmapsto f^*(\mathcal F)$, từ phạm trù các Môđun tự do địa
phương kiểu hữu hạn trên $S$ đến phạm trù các Môđun tự do địa phương kiểu hữu
hạn trên $X$, là đầy đủ và trung thành, và ảnh cốt yếu gồm các Môđun
$\mathcal G$ trên $X$ sao cho $f_*(\mathcal G)$ tự do địa phương và đồng cấu
chính tắc
\[
f^*\bigl(f_*(\mathcal G)\bigr)\longrightarrow\mathcal G
\]
là một đẳng cấu. Khẳng định thứ hai đã được chứng minh bằng lý thuyết hạ xuống
trong [4], I, B, 4.

Các kết quả của 2.2 cũng có thể được phát biểu như sau:

\result{HỆ QUẢ}{2.3} Dưới các điều kiện của 2.2, đồng cấu chính tắc
(1.9) $\operatorname{Pic}'_{X/S}\longrightarrow\operatorname{Pic}_{X/S}$ là
đơn ánh, và thậm chí song ánh nếu $X$ có một tiết diện địa phương trên $S$.
(Do đó, trong trường hợp sau, nhóm Picard tương đối
$\operatorname{Pic}(X/S)$ và nhóm Picard tương đối hạn chế
$\operatorname{Pic}'(X/S)$ đồng nhất với nhau.)

Kết hợp với 2.1, ta thu được:

\result{HỆ QUẢ}{2.4} Dưới các điều kiện của 2.2, ta có một dãy khớp
\[
0\longrightarrow\operatorname{Pic}(S)\longrightarrow\operatorname{Pic}(X)
\longrightarrow\operatorname{Pic}(X/S).
\]
Khi $X$ có một tiết diện trên $S$, đồng cấu cuối là toàn ánh, nghĩa là khi ấy
ta có một đẳng cấu
\[
\operatorname{Pic}(X/S)\xrightarrow{\sim}
\operatorname{Pic}(X)/\operatorname{Pic}(S).
\]

\result{NHẬN XÉT}{2.5} Cho $f:X\longrightarrow S$ là một cấu xạ sao cho
$\mathcal O_S\xrightarrow{\sim}f_*(\mathcal O_X)$, và cho $g$ là một tiết
diện của $X$ trên $S$. Cho $\mathcal L$ là một Môđun khả nghịch trên $X$; gọi
$g$-rigid hóa của $\mathcal L$ một đẳng cấu
$\mathcal O_S\xrightarrow{\sim}g^*(\mathcal L)$, và gọi
\emph{Môđun khả nghịch được $g$-rigid hóa} một Môđun khả nghịch $\mathcal L$
trên $X$ được trang bị một phép $g$-rigid hóa. Mọi tự đẳng cấu của một cấu
trúc như vậy đều tầm thường, và
$\operatorname{Pic}'(X/S)$ đồng nhất với tập các lớp đẳng cấu của những Môđun
khả nghịch được $g$-rigid hóa trên $X$\footnote{Bản in ghi «trên $S$»; phép
rigid hóa bởi một tiết diện $g:S\to X$ ở đây liên quan đến Môđun
$\mathcal L$ trên $X$.}. (Chính sự kiện này đã cho phép dùng lý thuyết hạ
xuống để chứng minh 2.2(ii).) Điều này cho một cách hiểu mới về
$\operatorname{Pic}(X/S)$, ít nhất khi $f$ còn tựa compact và tách, do đó
$\operatorname{Pic}(X/S)\simeq\operatorname{Pic}'(X/S)$ theo 2.3.
\src{148}
\result{NHẬN XÉT}{2.6} Cho $f:X\longrightarrow S$ là một cấu xạ như trong
2.2, và cho $S'\longrightarrow S$ là một cấu xạ phẳng trung thành và tựa
compact sao cho tồn tại một $S$-cấu xạ $S'\longrightarrow X$, nghĩa là sao cho
tồn tại một tiết diện của $X'=X\times_SS'$ trên $S'$. Đặt
$S''=S'\times_SS'$, $X''=X\times_SS''$, và xét dãy khớp
\[
\operatorname{Pic}(X/S)\longrightarrow\operatorname{Pic}(X'/S')
\rightrightarrows\operatorname{Pic}(X''/S'').
\]
Áp dụng 2.4, ta thu được một dãy khớp
\[
\operatorname{Pic}(X/S)\longrightarrow
\operatorname{Pic}(X')/\operatorname{Pic}(S')
\rightrightarrows
\operatorname{Pic}(X'')/\operatorname{Pic}(S''),
\]
đặc biệt, mọi phần tử của nhóm Picard tương đối đã «đến từ» một phần tử của
$\operatorname{Pic}(X')$. Như vậy điều này đơn giản hóa đáng kể mô tả nhóm
Picard tương đối đã cho trong mục trước, và tương tự đối với hàm tử Picard của
$X$ trên $S$, vì với mọi $T$ trên $S$, ta có thể áp dụng điều trên cho
$X\times_ST/T$ và cấu xạ $T'=S'\times_ST\longrightarrow T$. Chẳng hạn, nếu
chính $f$ phẳng trung thành, ta có thể lấy $S'=X$; điều này cho phép, khi $f$
còn thuộc kiểu hữu hạn (tương ứng trơn, v.v.), trong mô tả hàm tử Picard tương
đối $T\longmapsto\operatorname{Pic}(X\times_ST/T)$, chỉ cần xét các phép đổi
cơ sở $T'\longrightarrow T$ thuộc kiểu hữu hạn (tương ứng trơn, v.v.). Khi
$f$ xạ ảnh và phẳng, $S$ Noether địa phương, ta chứng minh được rằng có thể
chọn $S'\longrightarrow S$ ở trên sao cho $S'$ là hợp rời của các phủ phẳng
$S'_i$ của các tập mở $S_i$ của $S$ phủ $S$; nếu $f$ còn khả phân, có thể lấy $S'_i$
étale trên $S_i$.

\fsection{3}{Định lý tồn tại chính: phát biểu}

Ta chưa có, ngay cả dưới dạng phỏng đoán, một phát biểu về sự tồn tại của các
tiền lược đồ Picard bao quát mọi trường hợp đã biết. Một điều kiện «hầu như
cần thiết», nếu có thể nói như vậy, là $f:X\longrightarrow S$ phải riêng (bảo
đảm các tính chất hữu hạn thiết yếu) và phẳng. Các điều kiện này không đủ,
ngay cả khi $S$ là phổ của đại số các số đối ngẫu $k[t]/(t^2)$ trên một trường
$k$ (chẳng hạn trường các số phức $\mathbf C$), và $X$ có chiều 1. Vào thời
điểm viết thuyết trình này, các định lý tồn tại quan trọng nhất cho tiền lược
đồ Picard được suy ra từ định lý sau:

\result{ĐỊNH LÝ}{3.1} Cho $f:X\longrightarrow S$ là một cấu xạ giữa các tiền
lược đồ Noether địa phương. Giả sử:

\noindent (i) $f$ xạ ảnh;
\src{149}
\noindent (ii) $f$ phẳng;

\noindent (iii) các thớ hình học của $f$ là nguyên.

Dưới các điều kiện này, $\operatorname{Pic}_{X/S}$ tồn tại.

Chứng minh, sẽ được phác thảo trong hai mục tiếp theo, đồng thời sẽ chỉ ra điều
sau: cho $\xi$ là tiết diện của $\operatorname{Pic}_{X/S}$ tương ứng với một
bó rất rộng $\mathcal O_X(1)$ trên $X/S$ (nghĩa là được cảm sinh bởi một phép
nhúng xạ ảnh $X\longrightarrow\mathbf P(\mathcal E)$); khi đó tồn tại một tập
con mở $U$ của $\operatorname{Pic}_{X/S}$, là hợp rời của các tập con mở tựa
xạ ảnh trên $S$, sao cho $U$ ổn định dưới phép tịnh tiến bởi $\xi$, và
$\operatorname{Pic}_{X/S}$ là hợp tăng của các tập con mở $U-n\xi$ (tất cả
đều đẳng cấu với $U$). Đặc biệt suy ra rằng dưới các điều kiện của 3.1,
$\operatorname{Pic}_{X/S}$ tách trên $S$.

\result{NHẬN XÉT}{3.2} Các ví dụ cho thấy (chẳng hạn với $S$ là phổ của một
vành định giá rời rạc, và $X$ có chiều tương đối 1 trên $S$) rằng nếu trong
3.1 ta bỏ giả thiết (iii), thay nó bằng giả thiết yếu hơn rằng với mọi
$s\in S$, đồng cấu
\[
k(s)\longrightarrow H^0(X_s,\mathcal O_{X_s})
\]
là một đẳng cấu, thì $\operatorname{Pic}_{X/S}$ không nhất thiết tách trên
$S$; điều này xảy ra cả trong trường hợp các thớ hình học của $f$ rút gọn,
nhưng một thớ tổng quát hình học nguyên «vỡ» do chuyên biệt hóa thành hai thành
phần bất khả quy, lẫn trong trường hợp các thớ hình học của $f$ bất khả quy,
nhưng một thớ tổng quát hình học nguyên chuyên biệt hóa thành một «thớ bội».
Trường hợp thứ nhất xảy ra chẳng hạn với một đường conic suy biến thành hai
đường thẳng cắt nhau; một ví dụ của trường hợp thứ hai do D. Mumford cung cấp,
với một đường cong elliptic suy biến thành một đường cong elliptic kép. Các ví
dụ này đúng trong mọi đặc số.

\result{NHẬN XÉT}{3.3} Dưới các điều kiện của 3.1, tôi không biết
$\operatorname{Pic}_{X/S}$ có phải là một hợp rời của các tập con mở thuộc
kiểu hữu hạn, do đó tựa xạ ảnh, trên $S$ hay không.\footnote{Bản đính chính
năm 1962 cho biết D. Mumford đã giải câu hỏi này theo hướng khẳng định, bằng
lý thuyết các lược đồ Picard được phát triển ở đây.} Lưu ý rằng việc xét các
đa thức Hilbert $Q\in\mathbf Q[t]$ cho phép, như trong trường hợp các lược đồ
Hilbert ([4], IV), phân tích $\operatorname{Pic}_{X/S}$ thành một hợp rời của
các tập con mở $\operatorname{Pic}^{Q}_{X/S}$, và có vẻ hợp lý rằng các tập
sau là tựa xạ ảnh trên $S$; ít nhất ta sẽ thấy điều đó trong thuyết trình kế
tiếp khi $f$ là một cấu xạ trơn. Cần lưu ý rằng nếu thay giả thiết (i) bằng
giả thiết: $X$ xạ ảnh địa phương trên $S$ (điều này đủ để 3.1 đúng, vì bài toán
tồn tại của $\operatorname{Pic}_{X/S}$ hiển nhiên có tính địa phương trên
$S$), thì ngược lại, dễ dàng
\src{150}
đưa ra những ví dụ trong đó $\operatorname{Pic}_{X/S}$ chứa các thành phần
liên thông không thuộc kiểu hữu hạn trên $S$. Chẳng hạn, cho $X_0$ là một đa
tạp đại số xạ ảnh không kỳ dị trên một trường đóng đại số $k$, được trang bị
một tự đẳng cấu $u$ và một phần tử $\xi$ trong nhóm Néron--Severi của $X_0$,
sao cho các $u^n(\xi)$ đôi một khác nhau. Chẳng hạn có thể lấy $X_0$ là tích
của một đường cong elliptic $E$ với chính nó, và lấy $u$ là tự đẳng cấu
\[
(x,y)\longmapsto(x,y+x)
\]
của $E\times E$. Cho $S$ là hợp của hai đường cong bất khả quy không kỳ dị,
cắt nhau tại hai điểm $a$ và $b$. Trên $S$ có một phủ chính liên thông $P$ với
nhóm $\mathbf Z$; dùng các tác động của $\mathbf Z$ trên $X_0$ được xác định
bởi $u$, ta suy ra một phân thớ liên kết trên $S$ với thớ $X_0$ (tầm thường
trên $S-a$ và $S-b$), và trong trường hợp riêng đang xét, thực ra là một lược
đồ Abel trên $S$. Dễ thấy rằng $\operatorname{Pic}_{X/S}$, cũng là phân thớ
liên kết với $P$ và với các tác động của $\mathbf Z$ trên
$\operatorname{Pic}_{X_0/k}$ qua $u$, chứa một thành phần liên thông đẳng cấu
với
\[
P\times\operatorname{Pic}^{0}_{X_0/k},
\]
(trong đó $\operatorname{Pic}^{0}$ chỉ thành phần liên thông của phần tử đơn
vị trong $\operatorname{Pic}$), và thành phần này không thuộc kiểu hữu hạn
trên $S$. [Những hiện tượng tương tự cũng xảy ra trong nhiều trường hợp tiền
lược đồ Picard không tách trên $S$, như đã xét trong 3.2.]

\fsection{4}{Các ước Cartier tương đối và các phân thớ xạ ảnh}

Ta sẽ chỉ cần dùng các ước dương, và trong phần còn lại của mục này sẽ bỏ qua
tính từ bổ sung ấy.

Cho $X$ là một tiền lược đồ. Một \emph{ước Cartier} trên $X$, hay để đơn giản
chỉ gọi là ước trên $X$, là một tiền lược đồ con đóng $D$ của $X$, được xác
định bởi một Iđêan $\mathcal J$ là một Môđun khả nghịch, nghĩa là được sinh địa
phương bởi một tiết diện không là ước của không của $\mathcal O_X$. Với $D$,
ta liên kết một Môđun khả nghịch, cụ thể là
\[
\mathcal L(D)=\mathcal J^{-1},
\]
và đơn ánh chính tắc $\mathcal J\longrightarrow\mathcal O_X$ cho một đồng cấu
chính tắc
\[
s_D:\mathcal O_X\longrightarrow\mathcal J^{-1}=\mathcal L(D),
\qquad\text{nghĩa là}\quad s_D\in\Gamma(X,\mathcal L(D)).
\]
Hơn nữa, dữ liệu của một ước về cơ bản tương đương với dữ liệu của một Môđun
khả nghịch $\mathcal L$ trên $X$, được trang bị một tiết diện $s$ không là ước
của không ở mọi nơi; ngược lại, với một cặp như vậy $(\mathcal L,s)$, ta liên
kết «ước» của $s$, ký hiệu là $\operatorname{div}(s)$. Với một
$\mathcal L$ khả nghịch đã cho trên $X$, tập các
\src{151}
ước $D$ xác định $\mathcal L$ tương ứng song ánh với tập thương
\[
\Gamma(X,\mathcal L)^*/\Gamma(X,\mathcal O_X^*),
\]
trong đó $\Gamma(X,\mathcal L)^*$ chỉ bộ phận của
$\Gamma(X,\mathcal L)$ gồm các tiết diện không là ước của không ở mọi nơi.

Bây giờ giả sử ta có một cấu xạ thuộc kiểu hữu hạn địa phương
$f:X\longrightarrow S$, và để đơn giản giả sử $S$ Noether địa phương. Cho
$\mathcal J$ là một Iđêan coherent trên $X$, $D$ là lược đồ con của $X$ do nó
xác định, $x\in X$, và $s=f(x)$. Ta chứng minh được rằng các điều kiện sau
tương đương:

\noindent (i) $\mathcal J$ khả nghịch tại $x$ (nghĩa là $\mathcal J_x$ được
sinh bởi một phần tử chính quy của $\mathcal O_{X,x}$) và $D$ phẳng trên $S$
tại $x$;

\noindent (ii) $X$ và $D$ phẳng trên $S$ tại $x$, và $D_s$ là một ước Cartier
trên thớ $X_s$ tại điểm $x$;

\noindent (iii) $X$ phẳng trên $S$ tại $x$, và $\mathcal J_x$ được sinh bởi
một phần tử $f_x$ cảm sinh trên $X_s$ một mầm không là ước của không.

Khi đó ta nói $D$ là một \emph{ước Cartier tương đối} trên $X/S$ (hay chỉ là
ước tương đối trên $X/S$) tại điểm đang xét. Từ (i), lưu ý rằng khi đó $D$
cũng là một ước tương đối tại các điểm lân cận của $x$; do đó nếu $X$ và $D$
phẳng trên $S$, còn $D$ riêng trên $S$, thì tập các $s\in S$ sao cho $D_s$ là
một ước Cartier trong $X_s$ (nghĩa là sao cho $D$ là một ước Cartier tương đối
tại các điểm của $X_s$) là một bộ phận mở của $S$. Mặt khác, định nghĩa trên
đã được đặt ra chính xác để khái niệm ước Cartier tương đối ổn định qua một
phép đổi cơ sở tùy ý $S'\longrightarrow S$. Bây giờ xét tập
$\operatorname{Div}(X/S)$ các ước tương đối trên $X/S$, rồi hàm tử phản biến
theo $S'$ biến thiên trên $S$ được định nghĩa bởi
\[
\operatorname{Div}_{X/S}(S')=
\operatorname{Div}(X\times_SS'/S').
\]

Giả sử $X$ phẳng và riêng trên $S$. Khi đó, theo đặc trưng (ii) của các ước
Cartier tương đối, $\operatorname{Div}_{X/S}$ có thể được xem là một hàm tử
con của hàm tử $\operatorname{Hilb}_{X/S}$ được định nghĩa trong [4], IV, và
cấu xạ bao hàm
\[
\operatorname{Div}_{X/S}\longrightarrow\operatorname{Hilb}_{X/S}
\]
«biểu diễn được bởi các phép nhúng mở» (xem [5], IV, 3.13) theo các nhận xét
trên. Dùng định lý tồn tại chính của [4], IV, ta thu được:
\src{152}
\result{MỆNH ĐỀ}{4.1} Giả sử $f:X\longrightarrow S$ xạ ảnh và phẳng. Khi đó
hàm tử $\operatorname{Div}_{X/S}$ biểu diễn được, và chính xác hơn, được biểu
diễn bởi một tập mở của $\operatorname{Hilb}_{X/S}$.

Dùng (đối với một bó $\mathcal O_X(1)$ rất rộng đã cho trên $X/S$) phân tích
chính tắc của $\operatorname{Hilb}_{X/S}$ thành tổng các tập mở
$\operatorname{Hilb}^{Q}_{X/S}$ tương ứng với các đa thức Hilbert
$Q\in\mathbf Q[t]$, ta cũng suy ra một phân tích tương tự
\[
\operatorname{Div}_{X/S}=
\coprod_{Q\in\mathbf Q[t]}\operatorname{Div}^{Q}_{X/S}
\]
thành tổng các tập mở rời nhau, tựa xạ ảnh trên $S$.

Mặt khác, dùng ánh xạ $D\longmapsto\mathcal L(D)$, ta thu được một đồng cấu có
tính hàm tử
\begin{equation}
\operatorname{Div}_{X/S}\longrightarrow\operatorname{Pic}_{X/S}
\tag{+}
\end{equation}
mà ta sẽ nghiên cứu; dưới những điều kiện khá tổng quát, nó sẽ tỏ ra biểu diễn
được tương đối ([5], IV, 3). Vì vậy, hãy xuất phát với một phần tử $\xi$ của
$\operatorname{Pic}_{X/S}(S')$, để đơn giản ký hiệu giả sử $S'=S$, và chứng
minh rằng hàm tử con tương ứng của $\operatorname{Div}_{X/S}$ biểu diễn được.
Trước hết, xét trường hợp $\xi$ được xác định bởi một Môđun khả nghịch
$\mathcal L$ trên $X$. Ta sẽ giả sử $X$ riêng và phẳng trên $S$, và các thớ
hình học của $X$ trên $S$ là nguyên; điều này cũng kéo theo ([2], III, đoạn 7)
rằng ta có $\mathcal O_S\xrightarrow{\sim}f_*(\mathcal O_X)$ và quan hệ này
vẫn đúng sau mọi phép đổi cơ sở $S'\longrightarrow S$. Khi đó, các ước Cartier
tương đối $D$ trên $X/S$ sao cho $\mathcal L(D)$ và $\mathcal L$ xác định cùng
một phần tử của $\operatorname{Pic}(X/S)=\operatorname{Pic}_{X/S}(S)$, nghĩa
là theo 2.4, sao cho $\mathcal L(D)$ và $\mathcal L$ đẳng cấu địa phương trên
$S$, tương ứng song ánh với các tiết diện của bó thương
$f_*(\mathcal L)^*/\mathcal O_S^*$. Sự tương ứng này tương thích với các phép
đổi cơ sở. Mặt khác, các xét đoán tổng quát kiểu «Künneth» ở chỗ đã dẫn (xem
thêm [6]) cho thấy tính riêng của $X/S$ và tính phẳng của $\mathcal L$ trên
$S$ kéo theo sự tồn tại của một Môđun coherent $\mathcal Q$ trên $S$, được
xác định sai khác một đẳng cấu duy nhất, và một đẳng cấu các bó
\[
f_*(\mathcal L)\xrightarrow{\sim}
\operatorname{Hom}_{\mathcal O_S}(\mathcal Q,\mathcal O_S),
\]
\footnote{Bản in đặt $\mathcal O_X$ ở chỉ số dưới của
$\operatorname{Hom}$; vì $\mathcal Q$ là một Môđun trên $S$, đối ngẫu của nó
được lấy trên $\mathcal O_S$.} hơn nữa, việc tạo thành $\mathcal Q$ tương
thích với các phép đổi cơ sở. Ở đây $f_*(\mathcal L)^*$ chỉ bó con các tập hợp
của $f_*(\mathcal L)$ mà các tiết diện trên $U$ là những tiết diện của
$\mathcal L$ trên $f^{-1}(U)$ xác định các ước Cartier tương đối trên
\src{153}
$f^{-1}(U)/U$, nghĩa là cảm sinh các tiết diện không là ước của không trên các
$X_s$ ($s\in U$). Dùng giả thiết rằng các thớ $X_s$ là nguyên, điều này đơn
giản có nghĩa là các tiết diện cảm sinh trên các thớ $X_s$ không đồng nhất
bằng không, hay nói theo các đồng cấu địa phương
$\mathcal Q\longrightarrow\mathcal O_S$, rằng các đồng cấu ấy là toàn ánh
(Nakayama). Điều này chứng tỏ tập các tiết diện của
$f_*(\mathcal L)^*/\mathcal O_S^*$ tương ứng song ánh với tập các Môđun thương
khả nghịch của $\mathcal Q$, hay cũng vậy, theo định nghĩa của phân thớ xạ ảnh
$\mathbf P(\mathcal Q)$ liên kết với Môđun coherent $\mathcal Q$ (xem [5], V,
2), với tập các tiết diện của $\mathbf P(\mathcal Q)$ trên $S$. Mô tả này
tương thích với việc tạo thành các ảnh ngược, và do đó ta thu được định lý sau.

\result{ĐỊNH LÝ}{4.3} Cho $f:X\longrightarrow S$ là một cấu xạ riêng và phẳng
có các thớ hình học nguyên, với $S$ Noether địa phương, và cho $\mathcal L$ là
một Môđun khả nghịch trên $X$. Với mọi $S'$ trên $S$, gọi $T(S')$ là tập các
ước tương đối $D$ trên $X\times_SS'/S'$ sao cho $\mathcal L(D)$ đẳng cấu với
$\mathcal L\otimes_{\mathcal O_S}\mathcal O_{S'}$ địa phương trên $S'$ (nghĩa
là sao cho $\mathcal L(D)$ và
$\mathcal L\otimes_{\mathcal O_S}\mathcal O_{S'}$ xác định cùng một phần tử
của $\operatorname{Pic}(X\times_SS'/S')$). Khi đó tồn tại một Môđun coherent
$\mathcal Q$ trên $S$, được xác định sai khác một đẳng cấu duy nhất, sao cho
hàm tử $T$ được biểu diễn bởi phân thớ xạ ảnh $\mathbf P(\mathcal Q)$.

\result{HỆ QUẢ}{4.4} Nếu còn giả sử $f$ xạ ảnh, thì đồng cấu có tính hàm tử
$\operatorname{Div}_{X/S}\longrightarrow\operatorname{Pic}_{X/S}$ biểu diễn
được bởi các cấu xạ xạ ảnh.

Thật vậy, khi $X$ có một tiết diện (tương ứng, có một tiết diện địa phương)
trên $S$, thì đồng cấu trên được biểu diễn bởi các phân thớ xạ ảnh (tương ứng,
bởi các phân thớ xạ ảnh địa phương), nhờ 4.3 và 2.1. Trong trường hợp $f$ tựa
xạ ảnh, ta dễ dàng quy về trường hợp trên bằng một phương pháp hạ xuống, dùng
các tựa tiết diện hữu hạn và phẳng của $X$ địa phương trên $S$.

\result{NHẬN XÉT}{4.5} Dưới các điều kiện của 4.3, Môđun $\mathcal Q$ trên
$S$ nói chung không tự do địa phương; điều này được biểu thị bởi sự kiện rằng
chiều của các thớ rút gọn của $\mathcal Q$, nghĩa là chiều của
$H^0(X_s,\mathcal L_s)$ khi $s\in S$ biến thiên, có thể nhảy. Cho một Môđun
coherent $\mathcal Q$ trên tiền lược đồ Noether địa phương $S$, dễ kiểm tra
rằng với một $s\in S$ đã cho, $\mathcal Q$ tự do tại $s$ khi và chỉ khi
$\mathbf P(\mathcal Q)$ phẳng trên $S$ tại các điểm nằm trên $s$ (trong trường
hợp đó, nó thậm chí trơn trên $S$ tại các điểm nằm trên $s$); trong tình huống
đang xét, vì $\mathcal Q$ được định nghĩa theo $\mathcal L$ như trên, điều này
cũng có nghĩa là việc tạo thành ảnh trực tiếp $f_*(\mathcal L)$ «giao hoán với
phép đổi cơ sở trong
\src{154}
một lân cận của $s$», hay cũng vậy,
$f_*(\mathcal L)_s\longrightarrow H^0(X_s,\mathcal L_s)$ là toàn ánh. Chẳng
hạn điều này xảy ra nếu $H^1(X_s,\mathcal L_s)=0$. Với điều kiện các tiền lược
đồ đang xét tồn tại, các tiêu chuẩn này đặc biệt áp dụng cho tình huống phổ
dụng $\operatorname{Div}_{X/S}\longrightarrow\operatorname{Pic}_{X/S}$, và
cho một điều kiện cần và đủ, tương ứng một điều kiện đủ, để cấu xạ này trơn tại
một điểm đã cho của $\operatorname{Div}_{X/S}$.

\fsection{5}{Chứng minh định lý tồn tại chính}

Dưới các điều kiện của 3.1, chọn một Môđun $\mathcal O_X(1)$ rất rộng trên
$X/S$, và gọi $\xi$ là phần tử tương ứng của $\operatorname{Pic}(X/S)$. Để
viết gọn, đặt
\[
P(S')=\operatorname{Pic}(X\times_SS'/S'),
\]
và trước hết, để đơn giản, giả sử $X/S$ có một tiết diện. Gọi $P^+(S')$ là bộ
phận của $P(S')$ gồm các lớp của những $\mathcal L$ khả nghịch trên
$X\times_SS'$ sao cho
\[
R^if'_*(\mathcal L(n))=0\quad\text{với $i>0$ và mọi $n\geq0$},
\qquad
f'_*(\mathcal L(n))\neq0\quad\text{với mọi $n\geq0$}.
\]
Đó là những điều kiện ổn định qua đổi cơ sở, nên xác định một hàm tử con
$P^+$ của $P$, hiển nhiên ổn định dưới phép tịnh tiến bởi $\xi$. Dùng «các
định lý A và B» của Serre ([2], III, 2) và các kết quả tổng quát ([5], IV, 5),
dễ thấy rằng $P$ biểu diễn được khi và chỉ khi $P^+$ biểu diễn được, và khi ấy
$P^+$ sẽ được biểu diễn bởi một tập mở $U$ của tiền lược đồ
$\operatorname{Pic}_{X/S}$ biểu diễn $P$, còn tiền lược đồ này sẽ là hợp tăng
của các tập mở $U-n\xi$.

Để viết gọn, đặt $D=\operatorname{Div}_{X/S}$, và gọi $D^+$ là ảnh ngược của
$P^+$ qua cấu xạ chính tắc $D\longrightarrow P$. Do đó ta có một cấu xạ
\begin{equation}
D^+\longrightarrow P^+,
\tag{+}
\end{equation}
và ta đã biết $D^+$ được biểu diễn bởi một tập mở $D^+$ của tiền lược đồ
$D=\operatorname{Div}_{X/S}$ (sự tồn tại của nó được bảo đảm bởi 4.1). Từ 4.3
dễ suy ra cấu xạ (+) được biểu diễn bởi các cấu xạ trơn, xạ ảnh và toàn ánh
(chính xác hơn, bởi các phân thớ xạ ảnh liên kết với các Môđun tự do địa
phương không triệt tiêu ở mọi nơi): điều này là do khi $\mathcal L$ trên
$X\times_SS'$ như ở đầu mục này, thì $f'_*(\mathcal L)$ là một Môđun tự do
địa phương $\neq0$, mà việc tạo thành nó giao hoán với đổi cơ sở; với các ký
hiệu của 4.3, khi đó $\mathcal Q$ đẳng cấu với đối ngẫu của
$f'_*(\mathcal L)$. Dùng tiêu chuẩn tổng quát ([5], IV, 4.7), ta sẽ suy ra
tính biểu diễn được của $P^+$. Ở chỗ đã dẫn, ta sẽ lấy $\mathfrak C$ là lớp
các cấu xạ tiền lược đồ phẳng trung thành
\src{155}
và tựa compact (quả thật là các toàn cấu hiệu dụng theo [4], I B). Điều kiện
(a) ở chỗ đã dẫn, tức là (+) biểu diễn được bởi các cấu xạ là phần tử của
$\mathfrak C$, đã được kiểm chứng như vừa thấy; điều kiện (b) cũng vậy, nó có
nghĩa là hàm tử $P^+$ tương thích với hạ phẳng trung thành tựa compact, điều
này là hiển nhiên. Còn phải chứng minh điều kiện (c) ở chỗ đã dẫn, tức là quan
hệ tương đương $R$ trong tiền lược đồ $D^+$ suy ra từ cấu xạ (+) biểu diễn được
bởi $\mathfrak C$ là $\mathfrak C$-hiệu dụng, nghĩa là hiệu dụng và sao cho
$D^+\longrightarrow D^+/R$ thuộc $\mathfrak C$. Để làm điều này, trước hết
lưu ý rằng các tập mở $D^{+Q}$ của $D^+$ tương ứng với các đa thức Hilbert ảo
khác nhau $Q\in\mathbf Q[t]$ ổn định dưới $R$ (vì các thớ của $R$ liên thông);
điều này quy ta về việc chứng minh rằng với mọi $Q$, quan hệ tương đương $R^Q$
cảm sinh trong $D^{+Q}$ là $\mathfrak C$-hiệu dụng. Nhưng giờ đây $D^{+Q}$ tựa
xạ ảnh, còn quan hệ tương đương $R^Q$ xạ ảnh và phẳng. Do đó ta ở trong các
điều kiện áp dụng của [4], III, 6.1, và kết quả cần chứng minh suy ra.

Trong trường hợp tổng quát, khi $X/S$ không nhất thiết có một tiết diện, ta dễ
dàng quy về trường hợp trên bằng kỹ thuật hạ xuống, hoặc lặp lại chứng minh
trên với sửa đổi cần thiết trong định nghĩa của $P^+$.

\result{CÁC NHẬN XÉT}{5.1} Phương pháp được dùng về cơ bản là phương pháp của
Matsusaka để xây dựng xạ ảnh các đa tạp Picard. Kết quả được viện dẫn từ [4],
III về khả năng chuyển sang thương hiệu dụng cũng có thể dễ dàng suy ra từ
định lý tồn tại các lược đồ Hilbert (xem chẳng hạn [6]) (theo cách cổ điển, các
thương này được xây dựng bằng cách dùng các tọa độ Chow). Lưu ý rằng việc tạo
thành tập mở $\operatorname{Pic}^{+}_{X/S}$ của
$\operatorname{Pic}_{X/S}$ và sự phân tích nó thành các tập mở
$\operatorname{Pic}^{+Q}_{X/S}$ tựa xạ ảnh trên $S$, theo các đa thức Hilbert
của các ước xác định những Môđun khả nghịch đang xét, tương thích với đổi cơ
sở (điều này cho phép áp dụng kỹ thuật hạ xuống).

\result{CÁC NHẬN XÉT}{5.2} Không loại trừ khả năng
$\operatorname{Pic}_{X/S}$ tồn tại mỗi khi $f:X\longrightarrow S$ riêng và
phẳng, và các đồng cấu
\[
k(s)\longrightarrow H^0(X_s,\mathcal O_{X_s})\qquad(s\in S)
\]
là những đẳng cấu (điều kiện cuối này khi đó cũng có nghĩa là
$\mathcal O_S\xrightarrow{\sim}f_*(\mathcal O_X)$ và quan hệ này vẫn đúng sau
mọi phép đổi cơ sở $S'\longrightarrow S$). Ít nhất đây là điều ta chứng minh
được trong khuôn khổ các không gian giải tích khi $f$ còn xạ ảnh, bằng một
phương pháp khác (phương pháp của Chow, nếu không nhầm) được giải thích trong
[5], IX, 3.1.\footnote{Bản đính chính năm 1962 cho biết phỏng đoán tồn tại
được phát biểu ở đầu nhận xét này là sai, nhưng các phương pháp của Mumford cho
một định lý yếu hơn đôi chút.} Trong phương pháp này, việc chuyển sang
\src{156}
thương theo một quan hệ tương đương riêng và phẳng trong một tập mở của lược
đồ các ước được thay bằng việc lấy thương theo nhóm xạ ảnh trong lược đồ các
phép nhúng của $X$ vào $\mathbf P^r_S$. Phương pháp này có lẽ có thể thích nghi
với trường hợp các lược đồ, bằng cách dùng các kết quả của Mumford về việc lấy
thương theo nhóm xạ ảnh [8]; hiện nay chỉ có chứng minh thành văn khi $X$ có
«nhiều tiết diện địa phương» trên $S$, chẳng hạn khi $X$ khả ly trên một vành
địa phương đầy đủ có trường thặng dư đóng đại số.\footnote{Bản in bỏ cụm «có
trường thặng dư đóng đại số»; bản đính chính nói rõ rằng phản ví dụ của Mumford
khiến hạn chế này trở nên không thể thiếu.} Về nguyên tắc, phương pháp đang xét
có phạm vi tổng quát hơn, vì nó cũng cho sự tồn tại của các tiền lược đồ Picard
trong những trường hợp chúng không tách, và do đó phương pháp thứ nhất tất yếu
thất bại. (Về mặt kỹ thuật, khó khăn xuất phát từ sự kiện rằng khi các thớ hình
học của $f$ không còn nguyên, thì hàm tử được xét trong 4.3 không còn được biểu
diễn bởi chính phân thớ xạ ảnh $\mathbf P(\mathcal Q)$, mà bởi một tập mở của
nó; điều này dẫn đến bài toán tế nhị về việc lấy thương theo một quan hệ tương
đương phẳng nhưng không riêng.)

\result{NHẬN XÉT}{5.3} Lưu ý rằng chứng minh được đưa ra ở đây không viện đến
phép dựng trước các đa tạp Jacobi của các đường cong hay các họ đường cong,
cũng không viện đến lý thuyết các đa tạp Abel hay các lược đồ Abel, và do đó
khác về bản chất với những trình bày truyền thống như trong sách của Lang [7]
hay thuyết trình của Chevalley [1], vốn đi theo con đường do A. Weil phác thảo.
Ngay cả trong trường hợp các đa tạp Jacobi của những đường cong không kỳ dị
trên một trường đóng đại số (chẳng hạn trường các số phức), phép dựng đa tạp
Jacobi được đưa ra ở đây là phép dựng duy nhất đã biết cho những tính chất rất
mạnh mà ta đã lấy làm định nghĩa ở mục 1 (về cơ bản là các tính chất của
Chevalley, nhưng có tính đến các «đa tạp tham số» với phần tử lũy linh). Việc
phép dựng các lược đồ Picard phải đi trước chứ không đi sau lý thuyết các đa
tạp Abel là rõ ràng a priori, vì nói chung các lược đồ Picard không phải là các
lược đồ Abel cũng không quy được về chúng; điều này đã thấy trong trường hợp
các đường cong kỳ dị trên một trường đóng đại số, nơi xuất hiện các «đa tạp
Jacobi suy rộng» của Rosenlicht, vốn không phải là các đa tạp Abel. Hơn nữa, lý
thuyết các đa tạp Abel, và tổng quát hơn lý thuyết các lược đồ Abel, trở nên
đơn giản hơn nhiều khi đã có một lý thuyết tổng quát về các lược đồ Picard. Đặc
biệt, lý thuyết đối ngẫu cho các lược đồ Abel, nhất là các kết quả kiểu Cartier,
khi đó trở nên gần như hình thức (xem chẳng hạn [10]).
\src{157}
\result{CÁC NHẬN XÉT}{5.4} «Nguyên lý tương thích» của Igusa đối với đa tạp
Jacobi của một đường cong suy biến thành một đường cong kỳ dị chỉ có thể được
hiểu đúng như một định lý tồn tại của lược đồ Picard của một lược đồ đường cong
tương đối $X/S$, không nhất thiết trơn trên $S$. Vì vậy, đó là một trường hợp
riêng của định lý tồn tại chính 3.1 khi đường cong chuyên biệt là nguyên (nghĩa
là, theo ngôn ngữ cổ điển, bất khả quy với bội 1). Lưu ý rằng hiện nay trường
hợp một đường cong đặc biệt khả quy (ngay cả khi các thành phần có bội 1, nghĩa
là khi đường cong đặc biệt khả ly trên trường thặng dư) vẫn nằm ngoài các định
lý tồn tại đã biết, trừ trường hợp ta ở trên một vành định giá rời rạc đầy đủ
có trường thặng dư đóng đại số; so sánh 5.2. Bài toán tồn tại này chắc chắn sẽ
đặt ra khi xây dựng bằng hình học, trong lý thuyết lược đồ, các «phép compact
hóa» Baily--Satake của các lược đồ môđuli đường cong giống $g$. (Hiện nay phép
compact hóa này chỉ được biết với $g=1$, nhờ công trình của Igusa.)

\fsection{6}{Các định lý tồn tại tương đối}

Ở đây ta sẽ phác thảo một số trường hợp hữu ích trong đó sự tồn tại của một số
lược đồ Picard kéo theo sự tồn tại của một số lược đồ khác, nhờ đó có thể suy
ra từ định lý chính 3.1 nhiều định lý tồn tại khác.

\result{MỆNH ĐỀ}{6.1} Cho $f:X\longrightarrow S$ là một cấu xạ phẳng và xạ ảnh
sao cho trong phân tích Stein $f=f''f'$, cấu xạ
$f':X\longrightarrow S'$ phẳng và có các thớ hình học nguyên (do đó thỏa các
giả thiết của 3.1), còn cấu xạ hữu hạn $f'':S'\longrightarrow S$ phẳng. Khi đó
$\operatorname{Pic}_{X/S}$ tồn tại và (với các ký hiệu được đưa vào trong [4],
II, C 2) ta có một đẳng cấu chính tắc
\[
\operatorname{Pic}_{X/S}\xrightarrow{\sim}
\Pi_{S'/S}\operatorname{Pic}_{X/S'}.
\]

Để chứng minh, trước hết ta thiết lập một đẳng cấu các hàm tử
\[
\operatorname{Pic}_{X/S}\xrightarrow{\sim}
\Pi_{S'/S}\operatorname{Pic}_{X/S'},
\]
rồi dùng 3.1, vốn kéo theo rằng $\operatorname{Pic}_{X/S'}$ biểu diễn được, và
dùng cấu trúc đã nêu ở mục số\textsuperscript{} 3 của
$\operatorname{Pic}_{X/S'}$ (kéo theo rằng mọi bộ phận hữu hạn của một thớ của
$\operatorname{Pic}_{X/S'}$ trên $S$ đều nằm trong một tập mở afin) để có sự
tồn tại của $\Pi_{S'/S}\operatorname{Pic}_{X/S'}$.
\src{158}
Chẳng hạn, khi $X$ là một lược đồ tổng của các lược đồ $X_i$ trên $S$ thỏa các
điều kiện của 3.1, phát biểu 6.1 quy về phát biểu hiển nhiên
\[
\operatorname{Pic}_{X/S}\xrightarrow{\sim}
\prod_i\operatorname{Pic}_{X_i/S}.
\]

\result{HỆ QUẢ}{6.2} Cho $f:X\longrightarrow S$ là một cấu xạ xạ ảnh và phẳng
có các thớ hình học nguyên địa phương (chẳng hạn một cấu xạ xạ ảnh và chuẩn),
thì $\operatorname{Pic}_{X/S}$ tồn tại.

Hơn nữa, trong trường hợp này $S'$ là một phủ étale của $S$ (đúng mỗi khi $f$
khả ly, nghĩa là phẳng với các thớ hình học rút gọn), và ta kiểm tra được rằng
định lý cấu trúc phát biểu ở mục số\textsuperscript{} 3 cho
$\operatorname{Pic}_{X/S}$ vẫn đúng, nhờ cấu trúc tương tự của
$\operatorname{Pic}_{X/S'}$.

Mặt khác, một thủ tục hạ xuống cho một định lý tồn tại tương đối, mà phạm vi
của nó phụ thuộc vào lời giải cho các bài toán hạ không phẳng được nêu trong
[4], I, A (c), và ở đây ta chỉ trình bày tường minh một trường hợp riêng:

\result{MỆNH ĐỀ}{6.3} Cho $f:X\longrightarrow S$ là một cấu xạ riêng, và cho
$X_1$, $X_2$ là hai tiền lược đồ con của $X$ phẳng trên $S$, được xác định bởi
hai Iđêan coherent $\mathcal I_1$ và $\mathcal I_2$ sao cho
$\mathcal I_1\cap\mathcal I_2=(0)$ và
$\mathcal O_X/(\mathcal I_1+\mathcal I_2)$ phẳng trên $S$ (nghĩa là tiền lược
đồ con của $X$ vốn là hợp của $X_1$ và $X_2$ thì bằng $X$, còn giao $Z$ của chúng
phẳng trên $S$).\footnote{Chữ cái cuối được in là «$X$», nhưng điều kiện ngay
trước đó nói về tính phẳng của $\mathcal O_X/(\mathcal I_1+\mathcal I_2)$ trên
$S$ và do đó xác định $Z$ phẳng trên $S$.} Giả sử thêm rằng với mọi $s\in S$,
các đồng cấu
\[
k(s)\longrightarrow H^0(X_{i,s},\mathcal O_{X_{i,s}}),\qquad i=1,2,
\]
là song ánh. Khi đó đồng cấu tự nhiên của các hàm tử
\[
\operatorname{Pic}_{X/S}\longrightarrow
\operatorname{Pic}_{X_1/S}\times\operatorname{Pic}_{X_2/S}
\]
biểu diễn được bởi các cấu xạ afin; do đó, nếu
$\operatorname{Pic}_{X_1/S}$ và $\operatorname{Pic}_{X_2/S}$ tồn tại, thì
$\operatorname{Pic}_{X/S}$ cũng tồn tại, và cấu xạ chính tắc
\[
\operatorname{Pic}_{X/S}\longrightarrow
\operatorname{Pic}_{X_1/S}\times\operatorname{Pic}_{X_2/S}
\]
là afin.

Bằng hạ phẳng trung thành, ta quy về trường hợp $Z$ có một tiết diện trên $S$,
do đó xác định các tiết diện của $X$, $X_1$, $X_2$ trên $S$, cho phép loại bỏ
các tự đẳng cấu trong những cấu trúc đang xét như đã giải thích trong 2.5.
Chứng minh khi đó chỉ cần lưu ý rằng dữ liệu của một Môđun khả nghịch
\src{159}
«được rigid hóa» $\mathcal L$ trên $X$ tương đương với dữ liệu của một bộ ba
$(\mathcal L_1,\mathcal L_2,u)$, trong đó $\mathcal L_i$ là một Môđun khả
nghịch «được rigid hóa» trên $X_i$, còn $u$ là một đẳng cấu từ
$\mathcal L_1|Z$ đến $\mathcal L_2|Z$, tương thích với các phép rigid hóa. Chỉ
cần kiểm tra rằng với $\mathcal L_1$, $\mathcal L_2$ cố định, dữ liệu của $u$
được biểu thị bởi một tiết diện của một lược đồ thích hợp trên $S$, afin trên
$S$, điều này là dễ. Từ 6.3, ta dễ dàng suy ra:

\result{HỆ QUẢ}{6.4} Cho $X$ là một lược đồ riêng và khả ly trên một trường
$k$, và cho $X_i$ là các thành phần bất khả quy của $X$. Nếu các
$\operatorname{Pic}_{X_i/k}$ tồn tại, thì $\operatorname{Pic}_{X/k}$ cũng tồn
tại, và cấu xạ chính tắc
\[
\operatorname{Pic}_{X/k}\longrightarrow\prod_i\operatorname{Pic}_{X_i/k}
\]
là afin.

Kết hợp với 6.2, điều này chẳng hạn chứng tỏ sự tồn tại của
$\operatorname{Pic}_{X/k}$ mỗi khi $X$ là một lược đồ xạ ảnh và khả ly trên
một trường $k$. Khi $X$ không còn khả ly trên $k$, ta cũng có một kết quả quy
giản, dùng một lập luận của Oort [11]. Phương pháp cũng áp dụng trên một lược đồ
cơ sở tùy ý (trường hợp hữu ích chẳng hạn để chứng minh trong thuyết trình tiếp
theo kết quả hữu hạn đã thông báo ở 3.3). Để tránh một phát biểu quá kỹ thuật,
ta chỉ xét trường hợp cơ sở là một trường:

\result{MỆNH ĐỀ}{6.5} Cho $X$ là một lược đồ riêng trên một trường $k$, và cho
$X_0$ là một lược đồ con có cùng tập hợp nền (do đó được xác định bởi một Iđêan
lũy linh trên $X$). Khi đó cấu xạ có tính hàm tử
\[
\operatorname{Pic}_{X/k}\longrightarrow\operatorname{Pic}_{X_0/k}
\]
biểu diễn được bởi các cấu xạ afin. Đặc biệt, nếu
$\operatorname{Pic}_{X_0/k}$ tồn tại, thì $\operatorname{Pic}_{X/k}$ cũng tồn
tại, và cấu xạ
$\operatorname{Pic}_{X/k}\longrightarrow\operatorname{Pic}_{X_0/k}$ là afin.

Kết hợp với 6.4, ta dễ dàng kết luận:

\result{HỆ QUẢ}{6.6} Cho $X$ là một tiền lược đồ xạ ảnh trên một trường $k$,
thì $\operatorname{Pic}_{X/k}$ tồn tại.

\result{NHẬN XÉT}{6.6} Rất có cơ sở để tin rằng với mọi lược đồ $X$ riêng trên
một trường $k$, $\operatorname{Pic}_{X/k}$ tồn tại.\footnote{Bản đính chính năm
1962 cho biết Murre vừa giải quyết bài toán này theo hướng khẳng định.} Các kết quả trên
quy bài toán này về trường hợp $k$ đóng đại số và $X$ nguyên. Khi đó ta biết
tồn tại một lược đồ nguyên $X'$ xạ ảnh trên $k$ và một cấu xạ trội
$g:X'\longrightarrow X$ (bổ đề Chow). Do đó chỉ cần chứng minh rằng
\src{160}
cấu xạ có tính hàm tử tương ứng
\[
\operatorname{Pic}_{X/k}\longrightarrow\operatorname{Pic}_{X'/k}
\]
biểu diễn được (và hẳn là biểu diễn được bởi các cấu xạ afin), vì mặt khác ta
biết $\operatorname{Pic}_{X'/k}$ biểu diễn được. Điều này đặt ra những bài toán
hạ không phẳng hiện còn chưa giải quyết. Lưu ý rằng khi ta chỉ xét hạn chế của
hàm tử $\operatorname{Pic}_{X/k}$ lên các tiền lược đồ rút gọn (với $X$ riêng
và nguyên trên $k$ đóng đại số), ta quả thực thu được một hàm tử biểu diễn
được, như Chevalley [1] đã chứng minh khi $X$ chuẩn, và Seshadri [12] chứng minh
bằng một phương pháp hạ xuống trong trường hợp tổng quát. Nhưng với các ký hiệu
của ta, lược đồ do các tác giả ấy dựng không phải là
$\operatorname{Pic}_{X/k}$, mà là
$(\operatorname{Pic}_{X/k})_{\mathrm{réd}}$, lược đồ rút gọn tương ứng với
$\operatorname{Pic}_{X/k}$.

\result{NHẬN XÉT}{6.7} Tổng quát hơn, cho $f:X'\longrightarrow X$ là một cấu
xạ toàn ánh của các tiền lược đồ riêng trên $k$. Khi đó, các xét đoán về hạ
không phẳng dẫn đến phỏng đoán rằng
\[
\operatorname{Pic}_{X/k}\longrightarrow\operatorname{Pic}_{X'/k}
\]
là một cấu xạ afin,\footnote{Bản đính chính năm 1962 cho biết bài toán này có
lời giải khẳng định và dẫn chiếu đến đoạn cuối của phần Bình luận cho thuyết trình
tiếp theo.} điều này đặc biệt kéo theo (sau khi chia cho các thành phần liên
thông của các phần tử đơn vị) rằng đồng cấu tương ứng trên các nhóm
Néron--Severi là đơn ánh modulo xoắn. Đây là điều có thể kiểm tra bằng lý thuyết
giao khi $X$ và $X'$ không kỳ dị. Câu trả lời dường như chưa được biết trong
bất kỳ trường hợp nào khác.

\result{NHẬN XÉT}{6.8} Trái với điều người ta có thể nghĩ, việc xét các lược đồ
Picard của những lược đồ đại số có phần tử lũy linh là hữu ích, thậm chí không
thể thiếu, trong nhiều bài toán. Chẳng hạn, khi $X$ là một lược đồ xạ ảnh, giả
sử trơn trên $k$,\footnote{Bản in ghi «chính quy»; bản đính chính quy định «trơn
trên $k$».} và $Y$ là một tiết diện siêu phẳng, cần xét các «lân cận vô cùng
bé» mọi cấp $X_n$ của $Y$, và các lược đồ Picard
$\operatorname{Pic}_{X_n/k}$; khi $X$ bất khả quy có chiều $\geq4$ (tương ứng
$\geq3$), cấu xạ chính tắc
\[
\operatorname{Pic}_{X/k}\longrightarrow\operatorname{Pic}_{X_n/k}
\qquad(n\text{ lớn})
\]
là một đẳng cấu (tương ứng, cảm sinh một đẳng cấu trên các ảnh ngược của những
nhóm con xoắn của Néron--Severi), một kết quả sẽ hữu ích cho ta trong nghiên
cứu định tính các lược đồ Picard ở thuyết trình tiếp theo. Tương tự, việc xét
các lược đồ Picard của một số đường cong có phần tử lũy linh và các định lý nền
tảng của hình học hình thức [3] cho phép kiểm chứng, trong trường hợp đồng đặc
trưng, một phỏng đoán của Mumford, cụ thể là với mọi vành địa phương Noether
chuẩn đầy đủ $A$ có chiều 2, nhóm các lớp ước của $A$ có thể được xem là tập
các điểm hữu tỉ trên $k$ của một
\src{161}
nhóm đại số $G$ trên trường thặng dư $k$ (hơn nữa, $G$ được xác định chính tắc
khi ta chọn một trường hệ số trong $A$). Khi $A$ có chiều tùy ý, rất có cơ sở
để tin rằng tồn tại một nhóm pro-đại số trên $k$ đóng vai trò như $G$ ở trên;
trong trường hợp có thể «giải kỳ dị» $\operatorname{Spec}(A)$, nhóm này được
dựng như một giới hạn xạ ảnh của các lược đồ Picard của những lược đồ xạ ảnh
thích hợp (có phần tử lũy linh) trên $k$.

\section*{TÀI LIỆU THAM KHẢO}
\addcontentsline{toc}{section}{Tài liệu tham khảo}

\begin{description}
\item[{[1]}] CHEVALLEY (Claude). -- Sur la théorie de Picard,
\emph{Amer. J. of Math.}, tập 82, 1960, tr. 435--490.

\item[{[2]}] GROTHENDIECK (A.) và DIEUDONNÉ (J.). -- Éléments de géométrie
algébrique, Chapitre I et suivants. -- Paris, Presses universitaires de France,
1960--1961 (Institut des Hautes Études Scientifiques, Publications
mathématiques 4, 8, \ldots).

\item[{[3]}] GROTHENDIECK (Alexander). -- Géométrie formelle et géométrie
algébrique, \emph{Séminaire Bourbaki}, tập 11, 1958/59,
số\textsuperscript{} 182, 28 tr.

\item[{[4]}] GROTHENDIECK (Alexander). -- Techniques de descente et théorèmes
d'existence en géométrie algébrique, I--IV, \emph{Séminaire Bourbaki}, tập 12,
1959/60, số\textsuperscript{} 190, 29 tr. và số\textsuperscript{} 195, 22 tr.;
tập 13, 1960/61, số\textsuperscript{} 212, 20 tr. và
số\textsuperscript{} 221, 28 tr.

\item[{[5]}] GROTHENDIECK (Alexander). -- Techniques de construction en
géométrie analytique, I--X, \emph{Séminaire Cartan}, tập 13, 1960/61:
Déformation des structures analytiques complexes, số\textsuperscript{} 7--17.

\item[{[6]}] GROTHENDIECK (Alexander). -- Séminaire de géométrie algébrique,
do Lichtenbaum biên soạn, Harvard University, 1961 (sắp xuất bản).

\item[{[7]}] LANG (Serge). -- \emph{Abelian varieties}. -- New York,
Interscience Publishers, 1959 (Interscience Tracts in pure and applied
Mathematics, 7).

\item[{[8]}] MUMFORD (D.). -- An elementary theorem in geometric invariant
theory, \emph{Bull. Amer. math. Soc.}, tập 67, 1961, tr. 483--486.

\item[{[9]}] MUMFORD (D.). -- The topology of normal singularities of an
algebraic surface and a criterion for simplicity. -- Paris, Presses
universitaires de France, 1961 (Institut des Hautes Études Scientifiques,
Publications mathématiques, 9).

\item[{[10]}] MUMFORD (D.) và TATE (J.). -- Séminaire de géométrie algébrique,
Harvard University, học kỳ Xuân 1962 (sắp xuất bản).

\item[{[11]}] OORT (F.). -- Sur le schéma de Picard,
\emph{Bull. Soc. math. France}, tập 90, 1962 (sắp xuất bản).

\item[{[12]}] SESHADRI (C. S.). -- Luận án sắp xuất bản trong
\emph{Annali di Matematica}.
\end{description}

\par
